为了说明像~$\Th{Q}$这样的理论如何能够“谈论”可计算函数——尤其是关于可证性（通过哥德尔数）——我们证明了 $\Th{Q}$ \textbf{表示}所有可计算函数。所谓“$\Th{Q}$ 表示 $f(n)$”，指的是在语言~$\Lang{L_A}$ 中存在一个公式~$!A_f(x, y)$ 表达 $f(x) = y$，并且 $\Th{Q}$ 能够证明这一点。进而，这意味着只要 $f(n) = m$，就有 $\Th{Q} \Proves !A_f(\num{n},
\num{m})$ 和 $\Th{Q} \Proves \lforall[y][(!A_f(\num{n}, y) \lif y =
\num{m})]$。（这里，$\num{n}$ 是 $n$ 的\textbf{标准数词}，即后接 $n$ 个撇号 $\prime$ 的项 $\Obj 0^{\prime\dots\prime}$。项~$\num{n}$ 在标准模型~$\Struct{N}$ 中指称数~$n$，因此它是在 $\Lang{L_A}$ 中表示数 $n$ 的一种便捷方式。）为了证明 $\Th{Q}$ 表示所有可计算函数，我们回到了可计算函数的刻画，即那些可以从 $\Zero$、$\Succ$ 以及投影函数出发，通过复合、原始递归和正则极小化定义出来的函数。虽然证明基本函数的可表示性，以及证明由可表示函数通过复合和正则极小化定义的函数仍是可表示的，都相对容易，但原始递归的情形则较为困难。我们证明了，如果允许引入少数几个额外的基本函数（即加法、乘法以及等词~$=$ 的特征函数），实际上就可以避免通过原始递归来进行定义。这需要用到\textbf{β函数}，它使我们能够以基本的方式处理数的序列，并且其定义无需使用原始递归。